\documentclass[preprint,12pt]{elsarticle}

\usepackage{amssymb} \usepackage{amsmath} \usepackage{lineno} \usepackage{subcaption} \usepackage{hyperref} \usepackage{cleveref} \usepackage{booktabs} \usepackage{tabularx} \usepackage{array} \usepackage{siunitx} \usepackage{xcolor}

% Absorb the few unbreakable long tokens (siunitx units, hyphenated
% compounds) that otherwise overhang the 390pt text block.
\emergencystretch=2em

% Panel labels are typeset by LaTeX rather than baked into the figure files,
% so they match the document font and stay referenceable. Kept identical to
% the main text so both documents label panels the same way.
\renewcommand\thesubfigure{\Alph{subfigure}}
\captionsetup[subfigure]{labelformat=simple, labelsep=space, labelfont=bf,
                         justification=centering, skip=2pt}

% With no abstract or keywords, retain only the closing frontmatter rule.
\makeatletter
\let\si@pprintMaketitle\pprintMaketitle
\renewcommand{\pprintMaketitle}{%
  \begingroup
  \let\si@hrule\hrule
  \def\hrule{\let\hrule\si@hrule}%
  \si@pprintMaketitle
  \endgroup
}
\makeatother

\journal{Spectrochimica Acta Part A: Molecular and Biomolecular Spectroscopy}

\begin{document}

\begin{frontmatter}

\title{Supporting information for: \\ Participant-level Raman spectroscopy of circulating EV-enriched fractions reveals distributed spectral differences in ischemic stroke patients}


\end{frontmatter}


% --- Supplementary numbering (Figure S1, Table S1, section titles ...)
\renewcommand{\thefigure}{S\arabic{figure}}
\setcounter{figure}{0}
\renewcommand{\thetable}{S\arabic{table}}
\setcounter{table}{0}

\renewcommand{\thesection}{S\arabic{section}}
\renewcommand{\thesubsection}{S\arabic{section}.\arabic{subsection}}
\renewcommand{\thesubsubsection}{S\arabic{section}.\arabic{subsection}.\arabic{subsubsection}}
\setcounter{section}{0}


\section{Clinical outcome availability}

We compared stroke participants with observed and missing outcome data to assess potential selection bias in the clinical analyses (\cref{tab:cohortB}). Age comparisons used Mann--Whitney tests; sex and lesion-side comparisons used Fisher exact tests, as described in Section~2.2 of the main text. Missing baseline MoCA and $\Delta$MoCA data were more frequent among participants with left-sided lesions. These comparisons describe differences in outcome availability; they do not establish the cause of missingness.

\begin{table}[htbp]
\centering
\caption{Stroke participants with observed and missing data for each clinical endpoint. Age entries are medians; percentages use the observed or missing subgroup as the denominator. $p$-values are unadjusted two-sided Mann--Whitney $U$ tests for age and Fisher exact tests for sex and lesion side. MoCA: Montreal Cognitive Assessment; $\Delta$MoCA: 3-month minus baseline MoCA; RBANS: Repeatable Battery for the Assessment of Neuropsychological Status (3 months); ASPECTS: Alberta Stroke Program Early Computed Tomography Score.}
\label{tab:cohortB}

\begin{tabularx}{\textwidth}{@{}>{\raggedright\arraybackslash}p{0.16\textwidth}>{\centering\arraybackslash}p{0.07\textwidth}>{\centering\arraybackslash}p{0.07\textwidth}>{\centering\arraybackslash}X>{\centering\arraybackslash}X>{\centering\arraybackslash}X@{}}
\toprule
Endpoint & $n$ obs. & $n$ miss. & Age, obs./miss.\ ($p$) & \% female, obs./miss.\ ($p$) & \% left lesion, obs./miss.\ ($p$) \\
\midrule
Baseline MoCA & 72 & 25 & 66/68 (0.225) & 47/44 (0.820) & 31/76 ($<$0.001) \\
3-month MoCA & 92 & 5 & 67/66 (0.683) & 48/20 (0.369) & 40/80 (0.159) \\
$\Delta$MoCA & 71 & 26 & 66/68 (0.347) & 48/42 (0.653) & 31/73 ($<$0.001) \\
RBANS total & 62 & 35 & 67/66 (0.581) & 42/54 (0.291) & 40/46 (0.671) \\
ASPECTS & 87 & 10 & 67/70 (0.934) & 46/50 (1.000) & 45/20 (0.183) \\
\bottomrule
\end{tabularx}
\end{table}


\section{Loading-interval definition for the participant-level PCA}
\label{sec:loading_intervals}

We examined the first three principal components. Candidate loading intervals used the 90th percentile of absolute loading magnitude; intervals separated by less than 6~\si{\per\centi\metre} were merged, and intervals narrower than 8~\si{\per\centi\metre} were excluded. Published Raman assignments provided tentative biochemical annotations \cite{kaminska2022,imanbekova2021,czamara2015}. We repeated the analysis within the 97-participant stroke subset as a sensitivity check on the same cohort.


\section{Classification validation and permutation testing}

Scaling, hyperparameter selection, and probability-threshold determination were restricted to training data within a nested cross-validation framework (five outer folds repeated ten times; four inner folds). Model performance was evaluated using balanced accuracy, sensitivity, specificity, Matthews correlation coefficient (MCC), and ROC--AUC. Percentile 95\% intervals came from 1,000 participant-level out-of-bag bootstrap fits, with each draw trained on the in-bag sample and evaluated on participants excluded from that draw.

Permutation tests used a fixed five-fold partition for both the observed and permuted statistics. The primary permutation distribution was generated by permuting case-control labels within sex strata, while unrestricted permutations were used as a secondary approach. Sequential Monte Carlo testing \cite{besag1991} was conducted with a maximum of 1,000 permutations, a minimum of 200, and a 99\% Clopper--Pearson stopping interval. The actual number of draws per model is reported in \cref{tab:cvd_ml}. Reported $p$-values correspond to the locked-partition statistic, while the primary performance estimate is the mean across repeated nested folds. Benjamini--Hochberg correction was applied across the six model--representation combinations.

We assessed potential confounding by sex using a female-only refit of the primary PLS-DA model and a sex-only logistic regression baseline. These checks assess discrimination within one sex and discrimination based on sex alone; they do not establish a sex-independent disease effect. Age adjustment was unavailable because 49 of 50 controls lacked age data in the linked analysis dataset. A separate geometric diagnostic compared whole-spectrum PLS discriminant directions for sex and disease status using cosine similarity and residualized projections (Section~S6.2). This diagnostic used a different component-selection procedure from the primary nested classifier.


\section{Supplementary clinical association and prediction analyses}

We calculated Spearman correlations at each of the 1,802 normalized spectral channels for baseline MoCA, 3-month MoCA, $\Delta$MoCA, 3-month RBANS, ASPECTS, and age. Benjamini--Hochberg false-discovery-rate (FDR) correction was applied separately to each endpoint--spectral-block family. Maximum cluster-mass permutation tests used contiguous channels with two-sided channel-wise $p<0.05$ and 1,000 endpoint-label permutations. We required both channel-wise FDR and cluster-level support before reporting a supported region; uncorrected fallback candidate intervals served only as exploratory output. Holm adjustment controlled multiplicity across phase-level summaries.

Partial least squares (PLS) regression predicted baseline MoCA, 3-month MoCA, and $\Delta$MoCA from the fused spectrum. Component counts from 1 to 5 were selected within the training folds of five-fold nested cross-validation. Performance measures were cross-validated $R^2$, root mean squared error (RMSE), and mean absolute error (MAE). Endpoint-label permutation tests used adaptive stopping (minimum 200, maximum 1,000 draws), with Holm adjustment across the three endpoints. Whole-spectrum predictive performance was a prerequisite for interpreting sparse or minimal-panel models; none of the three endpoints met this criterion.

We defined a separate 10-feature tentatively assigned spectral panel after reviewing the whole-spectrum results. This panel was post hoc, not preregistered. Its designated association was the age- and sex-adjusted partial Spearman correlation between the tentatively assigned carotenoid-to-carbonyl spectral ratio (1525/1721~\si{\per\centi\metre}) and $\Delta$MoCA, with a participant-level bootstrap interval and permutation test. These analyses were exploratory.


\section{Supplementary software and reproducibility information}
Analyses used Python~3.12 with NumPy, SciPy, scikit-learn, pandas, statsmodels, and pybaselines under locked dependencies. Frozen configuration files specified preprocessing, participant-level PCA, the band screen, classification, and within-stroke analyses. Downstream analyses read preprocessed spectra from an analysis store identified by hashes of its configuration and input inventory. Result directories retained configuration and data-manifest snapshots, and reported figures were regenerated from saved participant-level tables.

Stochastic procedures used fixed seeds, including bootstrap intervals (2,000 resamples for rank-biserial effects and 1,000 out-of-bag fits for classification metrics), cross-validation folds, the 500-draw downsampling analysis in \cref{sec:measurement_imbalance}, and permutation tests. Tests of a single observed statistic used sequential Monte Carlo sampling \cite{besag1991} with a 99\% Clopper--Pearson stopping interval and 200--1,000 draws; \cref{tab:cvd_ml} reports the actual counts. The cluster-mass test in Section~S4 used a shared distribution of maximum cluster statistics across 1,000 permutations without early stopping. Access to processed participant-level matrices, frozen configurations, and analysis scripts follows the Data Availability statement in the main text; raw Raman and linked clinical data remain restricted.


\section{Classification model validation}
\subsection{Classifier performance}

The primary whole-spectrum PLS-DA model achieved a balanced accuracy of 0.655 (95\% CI 0.538--0.771), ROC--AUC of 0.711 (95\% CI 0.550--0.823), and MCC of 0.313 (\cref{tab:cvd_ml}). The ROC--AUC interval spans values close to chance and values above 0.8, indicating uncertainty in the magnitude of discrimination. The balanced accuracy on the locked partition used for permutation testing was 0.656, with sex-stratified $p=0.005$ ($q=0.0075$). This test evaluates the locked-partition statistic; it does not test the repeated-cross-validation mean directly. Random forest yielded higher point estimates (balanced accuracy 0.729; ROC--AUC 0.818), but no paired comparison established superiority. The 21-window models used the same participants and provided no independent validation.

\begin{table}[!t]
\centering
\caption{Exploratory case-control classification. Intervals are participant-level out-of-bag bootstrap 95\% intervals (1,000 resamples); $p$-values used sex-stratified permutation on one locked five-fold partition, with $q$ covering the six models. $B$ is the number of permutation draws actually taken.}
\label{tab:cvd_ml}
\scriptsize
\renewcommand{\arraystretch}{1.2}
\begin{tabularx}{\textwidth}{@{}>{\raggedright\arraybackslash}p{0.16\textwidth} >{\raggedright\arraybackslash}p{0.13\textwidth} >{\centering\arraybackslash}X >{\centering\arraybackslash}X cccc@{}}
\toprule
Model & Features & Balanced accuracy & ROC--AUC & MCC & $p$ & $q$ & $B$ \\
\midrule
PLS-DA (primary) & Whole spectrum & 0.655 (0.538--0.771) & 0.711 (0.550--0.823) & 0.313 & 0.005 & 0.0075 & 200 \\
Random forest & Whole spectrum & 0.729 (0.586--0.835) & 0.818 (0.701--0.911) & 0.434 & 0.005 & 0.0075 & 200 \\
$L_1$ logistic & Whole spectrum & 0.626 (0.500--0.738) & 0.667 (0.482--0.782) & 0.256 & 0.066 & 0.0660 & 1000 \\
\addlinespace
PLS-DA & 21 windows & 0.660 (0.513--0.743) & 0.693 (0.521--0.801) & 0.330 & 0.010 & 0.0119 & 200 \\
Random forest & 21 windows & 0.679 (0.578--0.800) & 0.781 (0.664--0.877) & 0.347 & 0.005 & 0.0075 & 200 \\
$L_1$ logistic & 21 windows & 0.642 (0.529--0.762) & 0.690 (0.533--0.802) & 0.286 & 0.005 & 0.0075 & 200 \\
\bottomrule
\end{tabularx}
\end{table}



Five models stopped at the 200-draw minimum (\cref{tab:cvd_ml}). The whole-spectrum $L_1$-penalized logistic model reached the 1,000-draw maximum because the stopping interval did not resolve its position relative to the significance threshold.

For whole-spectrum $L_1$-penalized logistic regression, the unrestricted permutation $p$-value was 0.040 and the sex-stratified value was 0.066. Thus, its unadjusted result crossed the 0.05 threshold when permutations preserved the sex composition of each group. The other five models did not show this threshold crossing. This sensitivity to the permutation scheme warrants caution but does not quantify the contribution of sex to classification performance.

Figure~4 of the main text presents the whole-spectrum classifiers. Panel~A averages the ten repeated nested-cross-validation ROC curves at common false-positive-rate grid points. The shading around PLS-DA shows $\pm1.96$ standard errors across repeats; because repeats reuse participants, this band describes variation across splits and is not a participant-level confidence interval. Panel~B shows 200 sex-stratified permutation draws and marks the repeated-cross-validation balanced accuracy (0.655). The reported $p=0.005$ instead uses the locked-partition statistic (0.656). The plot is therefore a visual summary rather than an exact display of the tested statistic. Only PLS-DA has a shaded band, so the plot cannot support a comparison based on overlapping model confidence bands.


\subsection{Sex-confound analyses}

The female-only PLS-DA analysis (45 stroke participants, 40 controls) achieved a balanced accuracy of 0.645 and ROC--AUC of 0.712 (a male-only model was not fitted because only 10 male controls were available). These estimates indicate discrimination within female participants but do not exclude confounding by other variables. The sex-only baseline achieved a balanced accuracy of 0.668 and ROC--AUC of 0.634. The spectral model had a higher ROC--AUC but lower balanced accuracy than this baseline; neither metric establishes calibration.

The geometric diagnostic compared whole-spectrum PLS directions for disease status and sex in the same 147 participants (\cref{fig:confound_geometry}). Component counts were selected using five-fold cross-validation on the full cohort before calculating out-of-fold scores; selection was not nested within the evaluation folds. The sex and disease discriminants used separate stratified five-fold partitions with the same random seed, yielding ROC--AUCs of 0.569 and 0.700, respectively. Full-cohort fits produced coefficient directions with cosine similarity 0.40. Residualizing the disease projection against the sex projection gave an apparent ROC--AUC of 0.688, calculated from full-data projections rather than held-out predictions. It cannot be compared with the out-of-fold value of 0.700 as a validated change in discrimination. These descriptive calculations show partial overlap between the fitted directions but neither bound confounding nor establish a sex-independent spectral effect.


\begin{figure}[htbp]
\centering
\begin{subfigure}[t]{0.48\textwidth}
\centering
\includegraphics[width=\linewidth]{figures/sex_disease_confound_auc.pdf}
\subcaption{Discriminant ROC--AUC}
\label{fig:confound_auc}
\end{subfigure}%
\hspace{0.02\textwidth}%
\begin{subfigure}[t]{0.48\textwidth}
\centering
\includegraphics[width=\linewidth]{figures/sex_disease_confound_projections.pdf}
\subcaption{Participant projections}
\label{fig:confound_projections}
\end{subfigure}
\caption{Exploratory comparison of disease- and sex-discriminant PLS directions ($n=147$). \subref{fig:confound_auc} five-fold out-of-fold ROC--AUCs after full-cohort component selection, using a separate stratified partition for each target. \subref{fig:confound_projections} participant projections from full-cohort fits. The projections are descriptive and do not represent held-out predictions.}
\label{fig:confound_geometry}
\end{figure}

Sex imbalance and missing control ages limit interpretation of the case-control results. Recruitment source, vascular risk factors, medication use, sampling time, and analytical batch may also differ between groups. Acute injury, inflammation, and hospitalization could contribute to spectral differences, but the design cannot separate these effects from disease status or establish a brain-derived EV contribution. The analyses support separation of the sampled groups; they do not establish diagnostic performance, causation, or clinical utility.


\section{Clinical association and prediction results}


\subsection{Descriptive severity-based spectral comparisons}
\Cref{fig:difference_spectrum} compares mean spectra at the severe and mild extremes of the stroke-severity distribution and shows the signed channel-wise effect size ($\mathrm{sign}(\Delta)\times|\text{Cohen's }d|$). The mean traces overlap, and the effect-size trace changes sign across both spectral blocks. This visual comparison does not establish a severity-associated spectral pattern. We applied no significance threshold or multiplicity correction because the contrast was descriptive.

\begin{figure}[ht]
\centering
\begin{subfigure}[t]{0.48\textwidth}
\centering
\includegraphics[width=\linewidth]{figures/difference_spectrum_fingerprint_mean.pdf}
\subcaption{Fingerprint, group means}
\label{fig:difference_fp_mean}
\end{subfigure}%
\hspace{0.02\textwidth}%
\begin{subfigure}[t]{0.48\textwidth}
\centering
\includegraphics[width=\linewidth]{figures/difference_spectrum_ch_stretch_mean.pdf}
\subcaption{C--H stretch, group means}
\label{fig:difference_ch_mean}
\end{subfigure}

\vspace{0.8em}

\begin{subfigure}[t]{0.48\textwidth}
\centering
\includegraphics[width=\linewidth]{figures/difference_spectrum_fingerprint_effect.pdf}
\subcaption{Fingerprint, effect size}
\label{fig:difference_fp_effect}
\end{subfigure}%
\hspace{0.02\textwidth}%
\begin{subfigure}[t]{0.48\textwidth}
\centering
\includegraphics[width=\linewidth]{figures/difference_spectrum_ch_stretch_effect.pdf}
\subcaption{C--H stretch, effect size}
\label{fig:difference_ch_effect}
\end{subfigure}
\caption{Extreme-group descriptive contrast between severe and mild ischemic stroke participants. \subref{fig:difference_fp_mean} and \subref{fig:difference_ch_mean} give the group mean normalized spectra for the fingerprint and C--H stretch blocks; \subref{fig:difference_fp_effect} and \subref{fig:difference_ch_effect} give the signed effect size ($\mathrm{sign}(\Delta)\times|\text{Cohen's }d|$) per channel, uncorrected. Descriptive context, not a significance-tested finding.}
\label{fig:difference_spectrum}
\end{figure}


\subsection{PLS prediction of MoCA outcomes}

PLS regression did not support prediction of any of the three MoCA endpoints (\cref{tab:clinical_pls}). Cross-validated $R^2$ was negative for baseline MoCA ($-0.335$), 3-month MoCA ($-0.067$), and $\Delta$MoCA ($-0.403$), indicating worse held-out performance than an outcome-mean reference. Adaptive permutation tests stopped after 200 draws per endpoint; neither raw nor Holm-adjusted $p$-values supported prediction. These results do not establish the absence of an association, but they provide no evidence of useful prediction under the present analysis.

\begin{table}[hbp]
\centering
\caption{Within-stroke PLS regression for MoCA outcomes. Metrics are from participant-level nested cross-validation; $p_{\mathrm{perm}}$ used adaptive endpoint-label permutation and $p_{\mathrm{Holm}}$ adjusts the three endpoints.}
\label{tab:clinical_pls}
\small
\renewcommand{\arraystretch}{1.2}
\begin{tabular}{@{}lrrrrrr@{}}
\toprule
Endpoint & $n$ & Components & CV $R^2$ & RMSE & MAE & $p_{\mathrm{perm}}$ / $p_{\mathrm{Holm}}$ \\
\midrule
Baseline MoCA & 72 & 1 & $-0.335$ & 6.273 & 4.641 & 0.910 / 1.000 \\
3-month MoCA & 92 & 1 & $-0.067$ & 5.851 & 4.682 & 0.214 / 0.642 \\
$\Delta$MoCA & 71 & 1 & $-0.403$ & 4.963 & 3.881 & 0.960 / 1.000 \\
\bottomrule
\end{tabular}
\end{table}


\subsection{Exploratory spectral panel analysis}

The designated association between the tentatively assigned carotenoid-to-carbonyl spectral ratio (1525/1721~\si{\per\centi\metre}) and $\Delta$MoCA gave $n=71$, partial Spearman $\rho=0.005$, a 95\% bootstrap interval of $[-0.255,\,0.277]$, and permutation $p=0.965$. A separate descriptive comparison of recovery versus decline in $\Delta$MoCA used Cliff's delta for the 10 panel features; all bootstrap intervals included zero (\cref{fig:panel_effect_sizes}). This trajectory contrast differs from the severity contrast in \cref{fig:difference_spectrum}. The post hoc panel analyses provide no supported clinical association.
%
\begin{figure}[ht] 
\centering 
\includegraphics[width=0.8\textwidth,trim={0cm 0cm 0cm 0.37cm}, clip]{figures/panel_effect_sizes.pdf}
\caption{Cliff's delta for the 10 post hoc tentatively assigned spectral panel features in the extreme-trajectory contrast (recovery versus decline in $\Delta$MoCA), with participant-level bootstrap 95\% intervals. Every interval includes zero.}
\label{fig:panel_effect_sizes} 
\end{figure}
%


\subsection{Unsupervised spectral descriptors}

As exploratory checks using alternative spectral summaries, we computed non-negative matrix factorization (NMF; three components, fingerprint block, stroke cohort) and per-participant Shannon spectral entropy without using endpoint labels, then tested their clinical associations post hoc. Figure~5D of the main text shows the entropy distribution. The largest raw Spearman correlation among NMF component weights was $|\rho|=0.102$ ($\Delta$MoCA, $q=0.883$); the largest age- and sex-adjusted partial correlation was $\rho=0.139$ for 3-month MoCA ($p=0.191$). Spectral entropy and the companion Gini-dispersion descriptor reached a maximum $|\rho|=0.093$ (all $q>0.94$). No test in either family survived Benjamini--Hochberg correction. These findings agree with the lack of supported associations in Section~3.6 of the main text, but reuse the same participants and do not constitute independent replication or prove that clinical associations are absent.


\section{Assessment of measurement-count imbalance}
\label{sec:measurement_imbalance}

Controls generally contributed fewer mapped-area measurements per spectral block than stroke participants. Two-stage median aggregation gave each participant one spectrum per block, but unequal measurement counts could affect its precision. We therefore downsampled each participant's measurements to the cohort-wide minimum of two per block (one participant had an excluded area), recomputed participant medians, and refitted the three participant-level PCA models across 500 seeded draws. Canonical and mean resampled PC1 variance were 75.8\% and 74.2\% for the fingerprint block (95\% resampling range 70.6--77.6\%), 67.1\% and 68.6\% for C--H (65.8--71.0\%), and 62.5\% and 59.3\% for the whole spectrum (56.3--62.4\%). Cosine similarity between resampled and canonical PC1 loadings ranged from 0.992 to 1.000 across the models. The dominant PC1 pattern persisted after matching measurement counts; this check does not address effects on classification, univariate band comparisons, clinical associations, or PC2/PC3 ordering.

\bibliographystyle{elsarticle-num}
\bibliography{cas-refs}


\end{document} 